% Chapter 7: Identification in the Limit
% ex_learning = Profile B (load-bearing), bc_learning = Profile C (revelation),
% mind_change_characterization = Profile B, online_pac_gold_separation = Profile D
%
% Texture: recursion-theoretic. Proofs use diagonalization, not concentration
% or combinatorial trees. Infinite horizon, no sample complexity, ordinal-valued
% mind-change bounds. This chapter should feel fundamentally different from
% Ch5 (PAC) and Ch6 (Online).

\chapter{Identification in the Limit}
\label{ch:gold}

In 1967, seventeen years before Valiant introduced PAC learning, E.\ Mark Gold posed a question that remains the oldest open research programme in computational learning theory: which classes of recursive functions can a machine identify from examples?

Gold's framework differs from PAC and online learning not merely in its definitions but in its \emph{proof technique}. PAC theory rests on concentration inequalities: the probabilistic convergence of empirical risk to true risk. Online learning rests on combinatorial game trees: the Littlestone dimension measures the depth of a binary tree that the adversary can force. Gold-style identification rests on \emph{diagonalization}: the recursion-theoretic technique of defeating every candidate learner by constructing an adversarial input that forces failure. The proof method shapes the entire chapter.

Three features distinguish this paradigm from everything that precedes it in this book:

\begin{enumerate}[leftmargin=*,itemsep=3pt]
\item \textbf{Infinite horizon.} The learner receives an infinite stream of data and must eventually converge. There is no sample complexity bound, no $\varepsilon$, no $\delta$. The question is whether the learner converges, not how fast.

\item \textbf{Ordinal-valued complexity.} When we ask ``how many times does the learner change its mind?'' the answer is not an integer but a \emph{transfinite ordinal}. This is the first appearance of ordinals in learning theory.

\item \textbf{A lattice of success criteria.} PAC learning has one definition (modulo agnostic, realizable, improper variants). Gold-style identification has a rich hierarchy: $\FIN < \Ex < \BC$, with anomalous, monotonic, and vacillatory learning branching off. The hierarchy itself is mathematical content.
\end{enumerate}

\begin{historical}
Gold's 1967 paper~\cite{Gold1967} was preceded by his 1965 paper on limiting recursion~\cite{Gold1965}, which introduced the idea that a computation could ``converge in the limit'' even if it never halts in the classical sense. The 1967 paper applied this idea to language learning. It was the first mathematical theory of learning from examples, predating Valiant's PAC model by 17 years, Vapnik and Chervonenkis's foundational work on uniform convergence by 4 years, and Littlestone's online model by 21 years. The diagonalization technique Gold used to prove his impossibility theorem later became standard in computational complexity theory, but Gold's use predates most of those applications.
\end{historical}

% ================================================================
% SECTION 1: Gold's Question
% ================================================================

\section{Gold's Question}
\label{sec:gold-question}

Fix a countable domain. In Gold's original setting, this is the set of all strings over a finite alphabet $\Sigma$, and the objects to be learned are formal languages $L \subseteq \Sigma^*$. We work in this setting when it clarifies the recursion-theoretic content, and in the general setting of concept classes $\mathcal{C} \subseteq 2^\N$ when it does not matter.

The learner receives data sequentially, one datum at a time, forever. Two data presentation modes are standard:

\begin{defn}[Text and Informant]
\label{def:text-informant}
Let $L \subseteq \Sigma^*$ be a language.
\begin{itemize}[leftmargin=*,itemsep=2pt]
\item A \emph{text} for $L$ is an infinite sequence $t = (t_0, t_1, t_2, \ldots)$ with $\{t_i : i \in \N\} = L$. Every element of $L$ appears at least once; no element of $\Sigma^* \setminus L$ ever appears. The text may contain repetitions and need not follow any fixed ordering.
\item An \emph{informant} for $L$ is an infinite sequence of labelled pairs $((s_0, b_0), (s_1, b_1), \ldots)$ where $\{s_i : i \in \N\} = \Sigma^*$ and $b_i = 1$ if and only if $s_i \in L$. Every string is eventually presented together with its membership status.
\end{itemize}
\end{defn}

Text presents only positive examples; informant presents both positive and negative examples. This distinction matters enormously: Gold showed that informant is strictly more powerful than text. We focus on text, which is the harder and more interesting case.

After seeing the first $n$ data points $(t_0, \ldots, t_{n-1})$, the learner outputs a hypothesis $h_n$. The hypothesis is an \emph{index}, a natural number naming a program that computes a language. The learner has no deadline, no accuracy target, and no probability distribution on targets. It simply sees more and more of $L$ and must eventually guess correctly.

% ================================================================
% SECTION 2: Ex-Learning and Gold's Impossibility Theorem
% ================================================================

\section{Ex-Learning and Gold's Impossibility Theorem}
\label{sec:ex-gold}

\begin{defn}[Explanatory Learning ($\Ex$)]
\label{def:ex-learning}
A learner $M$ \emph{$\Ex$-identifies} a language $L$ from text if, for every text $t$ for $L$, there exists an index $e$ and a time $n_0$ such that $M(t_0, \ldots, t_n) = e$ for all $n \geq n_0$, and $W_e = L$ (where $W_e$ is the language computed by program $e$).

A class $\mathcal{L}$ of languages is \emph{$\Ex$-identifiable} if there exists a single learner $M$ that $\Ex$-identifies every $L \in \mathcal{L}$ from text.\formal{FLT_Proofs.Criterion.Gold.EXLearnable}
\end{defn}

The definition requires \emph{syntactic convergence}: the learner must eventually output the same index forever. It is not enough to output a sequence of indices that all happen to compute the correct language; the index itself must stabilize. (Relaxing this requirement gives BC-learning; see \Cref{sec:bc-learning}.)

\begin{example}[$\Ex$-identification of finite languages]
\label{ex:finite-languages}
Let $\mathcal{L}_{\mathrm{fin}}$ be the class of all finite subsets of $\N$. Here is an $\Ex$-learner for $\mathcal{L}_{\mathrm{fin}}$: at time $n$, output an index for $\{t_0, \ldots, t_n\}$. After all elements of $L$ have appeared (which happens at some finite time $n_0$, since $L$ is finite), the set $\{t_0, \ldots, t_n\}$ equals $L$ for all $n \geq n_0$. The hypothesis stabilizes.
\end{example}

This simple algorithm illustrates an asymmetry: the learner does not know \emph{when} it has converged. It has no way to announce ``I am done.'' It merely stabilizes, silently. This lack of a convergence signal is what makes Gold's impossibility theorem possible.

\begin{thm}[Gold's Impossibility Theorem {\cite{Gold1967}}]
\label{thm:gold}
Let $\mathcal{L}$ be any class of languages that contains all finite languages and at least one infinite language. Then $\mathcal{L}$ is not $\Ex$-identifiable from text.\formal{FLT_Proofs.Theorem.Gold.gold_theorem}
\end{thm}

This is the centerpiece of the chapter. The proof is a diagonalization argument: we defeat \emph{every} candidate learner by constructing a text that forces it to fail. The construction is adversarial, not probabilistic: there is no distribution, no $\varepsilon$-net, no covering number. The adversary simply watches the learner and manipulates the future of the stream to prevent convergence.

\begin{proof}
Let $M$ be any learner. We construct a text $t$ for some language $L \in \mathcal{L}$ such that $M$ fails to $\Ex$-identify $L$ on $t$.

We build $t$ in stages. Let $L_\infty \in \mathcal{L}$ be the infinite language that $\mathcal{L}$ contains by assumption. Enumerate $L_\infty = \{w_0, w_1, w_2, \ldots\}$.

\medskip
\noindent\textbf{Stage 0.} Present $w_0$ to $M$. The learner outputs some hypothesis $h_0 = M(w_0)$.

\medskip
\noindent\textbf{Stage $k \geq 1$.} At this point we have presented the finite sequence $(w_0, \ldots, w_{n_{k-1}})$ for some $n_{k-1}$, and $M$ has output hypothesis $h_{n_{k-1}}$. Let $F_k = \{w_0, \ldots, w_{n_{k-1}}\}$ be the set of strings presented so far.

Consider two cases:
\begin{itemize}[leftmargin=*]
\item \textbf{Case A:} $W_{h_{n_{k-1}}} = F_k$ (the current hypothesis computes exactly the finite set seen so far). Then present $w_{n_{k-1}+1}$ (the next element of $L_\infty$), enlarging $F_k$. Let $n_k = n_{k-1} + 1$. Proceed to stage $k+1$.

\item \textbf{Case B:} $W_{h_{n_{k-1}}} \neq F_k$. Then the current hypothesis is already wrong for the finite language $F_k$. We may define $L = F_k$: since $F_k$ is finite and $\mathcal{L}$ contains all finite languages, $F_k \in \mathcal{L}$. We extend the text by repeating elements of $F_k$ forever. The learner $M$ has output a hypothesis $h_{n_{k-1}}$ with $W_{h_{n_{k-1}}} \neq F_k = L$, and the text for $L$ can be arranged so that $M$ never converges to a correct index (we continue to the next stage rather than stopping, as shown below).
\end{itemize}

The key observation is the \emph{diagonalization}: we never let the learner rest.

If Case~A occurs at every stage, then every element of $L_\infty$ is eventually presented, so $t$ is a text for $L_\infty$. At each stage, $M$'s current hypothesis $h_{n_k}$ is checked: does $W_{h_{n_k}} = F_k$? If yes, we expand. But $F_k \subsetneq L_\infty$ for all $k$ (since $L_\infty$ is infinite), so $W_{h_{n_k}} = F_k \neq L_\infty$. Therefore $M$ outputs a hypothesis that is wrong at every stage; it cannot converge to a correct index for $L_\infty$.

If Case~B occurs at some stage $k$, then $M$ has already failed for $F_k$. We commit to $L = F_k$ and repeat elements of $F_k$ forever. But we can do more: before committing, we wait to see if $M$ changes its mind. If $M$ later outputs a hypothesis $h'$ with $W_{h'} = F_k$, we switch to Case~A and add the next element of $L_\infty$. This forces $M$ to fail for $L_\infty$ instead.

In either case, $M$ fails. Since $M$ was an arbitrary learner, no learner $\Ex$-identifies $\mathcal{L}$.
\end{proof}

The proof has a recursive structure that repays careful study. The adversary maintains a ``threat'' at every stage: either the current language is the finite set seen so far, or it is the infinite language $L_\infty$. Whenever the learner commits to one possibility, the adversary switches to the other. The learner is forced to change its mind infinitely often, and $\Ex$-identification requires that it eventually stop changing its mind.

\begin{remark}[Diagonalization vs.\ concentration]
\label{rmk:diag-vs-conc}
Compare this proof to the lower bound for PAC learning (\Cref{ch:pac}). The PAC lower bound constructs a \emph{distribution} that fools the learner with positive probability; it is probabilistic. Gold's proof constructs a \emph{text} that fools the learner with certainty; it is adversarial. The PAC proof uses counting (how many hypotheses can be distinguished by $m$ samples); Gold's proof uses the halting problem implicitly (the adversary must check whether $W_e = F$ for arbitrary $e$). These are entirely different mathematical worlds.
\end{remark}

\begin{historical}
The proof technique of Theorem~\ref{thm:gold} is a \emph{priority argument} in the sense of recursion theory, though a simple one. More sophisticated priority arguments appear in the study of learning with resource bounds (Case and Smith~\cite{CaseSmith1983}). The connection between Gold-style learning and recursion theory runs deep: Barzdin\v{s}~\cite{Barzdinsh1974} showed that the class of languages $\Ex$-identifiable from informant is exactly the class of $\Sigma^0_2$-definable families, establishing a precise link to the arithmetical hierarchy.
\end{historical}

% ================================================================
% SECTION 3: The Hierarchy: BC > Ex > FIN
% ================================================================

\section{The Identification Hierarchy}
\label{sec:hierarchy}

Gold's impossibility theorem says that $\Ex$-identification is limited. Two natural responses: strengthen the success criterion (demand faster convergence) or weaken it (demand less of the final hypothesis). Both directions are fruitful.

\subsection{Finite Identification}
\label{sec:fin-learning}

\begin{defn}[Finite Identification ($\FIN$)]
\label{def:fin-learning}
A learner $M$ \emph{$\FIN$-identifies} a language $L$ from text if, for every text $t$ for $L$, there exists a time $n_0$ such that $M$ outputs exactly one hypothesis: $M$ outputs ``?'' (no guess) for $n < n_0$ and outputs a single index $e$ at time $n_0$ with $W_e = L$, never changing its mind. A class $\mathcal{L}$ is \emph{$\FIN$-identifiable} if a single learner $\FIN$-identifies every $L \in \mathcal{L}$.\formal{FLT_Proofs.Criterion.Gold.FiniteLearnable}
\end{defn}

$\FIN$ is the strongest identification criterion: the learner must produce the correct answer in finitely many steps with no subsequent revision. The class of $\FIN$-identifiable families is a strict subset of $\Ex$-identifiable families; for instance, the class of all finite languages is $\Ex$-identifiable (\Cref{ex:finite-languages}) but not $\FIN$-identifiable, since the learner cannot know when the last element has been seen.

\subsection{Behaviorally Correct Learning}
\label{sec:bc-learning}

\begin{defn}[Behaviorally Correct Learning ($\BC$)]
\label{def:bc-learning}
A learner $M$ \emph{$\BC$-identifies} a language $L$ from text if, for every text $t$ for $L$, there exists a time $n_0$ such that $W_{M(t_0, \ldots, t_n)} = L$ for all $n \geq n_0$. That is, from time $n_0$ onward, every hypothesis the learner outputs is \emph{extensionally correct} (computes the right language), but the indices may keep changing.\formal{FLT_Proofs.Criterion.Gold.BCLearnable}
\end{defn}

The distinction between $\Ex$ and $\BC$ is subtle but consequential. $\Ex$ requires \emph{syntactic} convergence: the index stabilizes. $\BC$ requires only \emph{semantic} convergence: the computed language stabilizes, even if the learner keeps switching between different programs that all compute the same language. At first glance, this seems like a minor relaxation. It is not.

\begin{thm}[Case--Smith {\cite{CaseSmith1983}}]
\label{thm:bc-strictly-stronger}
$\BC$ is strictly more powerful than $\Ex$: there exists a class of languages that is $\BC$-identifiable but not $\Ex$-identifiable from text.
\end{thm}

\begin{proof}
We construct a class $\mathcal{L}$ that separates $\BC$ from $\Ex$.

For each total recursive function $f : \N \to \N$, define the language $L_f = \{\langle n, f(n)\rangle : n \in \N\}$, where $\langle \cdot, \cdot \rangle$ is a computable pairing function. Let $\mathcal{L} = \{L_f : f \text{ is total recursive}\}$.

\medskip\noindent\textbf{$\mathcal{L}$ is $\BC$-identifiable.} Given a text $t$ for some $L_f$, at time $n$ the learner has seen finitely many pairs $\langle n_i, m_i\rangle$. Define the partial function $g_n$ by $g_n(n_i) = m_i$ for all pairs seen so far, and $g_n(k) = 0$ for all other $k$. Output an index for the language $L_{g_n} = \{\langle k, g_n(k)\rangle : k \in \N\}$. Once all pairs $\langle k, f(k)\rangle$ for $k \leq K$ have appeared (for any $K$), the hypothesis computes $L_f$ correctly on $\{0, \ldots, K\}$. As $n \to \infty$, $g_n$ converges pointwise to $f$. Since $f$ is total, for every $k$ there exists a time after which $g_n(k) = f(k)$. The language computed by the hypothesis is eventually $L_f$, though the \emph{index} keeps changing as new pairs are absorbed. This is $\BC$-identification.

\medskip\noindent\textbf{$\mathcal{L}$ is not $\Ex$-identifiable.} Suppose for contradiction that a learner $M$ $\Ex$-identifies $\mathcal{L}$. We diagonalize against $M$. Define a total recursive function $f$ as follows.

Begin presenting pairs $\langle 0, 0\rangle, \langle 1, 0\rangle, \langle 2, 0\rangle, \ldots$ (corresponding to the zero function). Wait until $M$ converges to some index $e$. Since $M$ $\Ex$-identifies $\mathcal{L}$ and the zero function is total recursive, $M$ must converge. Now $W_e = L_{\mathbf{0}}$ (the language of the zero function).

Modify $f$: set $f(k) = 1$ for some large $k$ not yet presented. This changes the target to a different total recursive function, but the text presented so far is consistent with both targets. The learner $M$ has already committed to index $e$, which computes $L_{\mathbf{0}} \neq L_f$. By a careful effectivization of this argument (choosing $k$ computably based on $M$'s behavior), we construct $f \in \mathcal{L}$ on which $M$ fails. This contradicts $M$ $\Ex$-identifying all of $\mathcal{L}$.

The essential point is that $\BC$-identification does not require the \emph{index} to stabilize, only the \emph{extension}. The class $\mathcal{L}$ exploits this: the learner must continually update its program as new pairs arrive, and the programs keep changing, but they all eventually compute the same language. $\Ex$-identification cannot tolerate this perpetual updating of indices.
\end{proof}

\begin{figure}[ht]
\centering
\begin{tikzpicture}[
    criterion/.style={draw, thick, rounded corners, minimum width=2.2cm, minimum height=0.85cm,
                      font=\small\bfseries, align=center},
    branch/.style={draw, rounded corners, minimum width=2cm, minimum height=0.7cm,
                   font=\small, align=center},
    strict/.style={-{Stealth[length=2mm]}, thick, sepgreen!70!black},
    weak/.style={-{Stealth[length=2mm]}, thin, dashed, notegray},
    witlabel/.style={font=\tiny, text=edgeorange, fill=white, inner sep=1.5pt, align=center},
    hdl/.style={font=\tiny, fill=white, inner sep=1pt, align=center}
]
% --- spine (layered, x=0; bottom-to-top = increasing power) ---
\node[criterion, fill=defblue!15] (fin) at (0,0)   {$\FIN$};
\node[criterion, fill=defblue!15] (ex)  at (0,2.7) {$\Ex$};
\node[criterion, fill=defblue!15] (bc)  at (0,5.4) {$\BC$};
% --- branches: restricts LEFT, extends RIGHT ---
\node[branch, fill=edgeorange!8] (mono)  at (-3.5,2.7) {Monotonic\\$\mathrm{Mon}\Ex$};
\node[branch, fill=edgeorange!8] (anom)  at ( 3.6,1.9) {Anomalous\\$\Ex^*$};
\node[branch, fill=goldcolor!10] (vac)   at ( 3.6,4.0) {Vacillatory\\$\mathrm{Vex}$};
\node[branch, fill=goldcolor!10] (trial) at ( 3.6,6.0) {Trial \&\\Error};
% --- spine edges: FIN<Ex SOLID (verified), Ex<BC DASHED (literature) ---
\draw[strict] (fin) -- node[hdl, left, xshift=-2pt]
   {all finite langs\\are $\Ex$, not $\FIN$\\{\tiny\leanname{ex_strictly_stronger_finite}}} (ex);
\draw[weak] (ex) -- node[hdl, left, xshift=-2pt]
   {Case--Smith 1983\\(\cref{thm:bc-strictly-stronger})\\\itshape literature, not kernel} (bc);
% --- relaxation edges: all DASHED, oriented by constraint-toggle direction ---
\draw[weak] (anom) -- node[witlabel, above, pos=0.45]{\rel{extends}\\$-$\,zero-error} (ex);
\draw[weak] (mono) -- node[witlabel, above]{\rel{restricts}\\$+$\,no-retraction} (ex);
\draw[weak] (vac) -- node[witlabel, below right, xshift=-3pt]{\rel{extends}} (ex);
\draw[weak] (vac) -- node[witlabel, above right, xshift=-3pt]{\rel{restricts}} (bc);
\draw[weak] (bc) -- node[witlabel, above right, xshift=-3pt]{$\Delta^0_2$\\(Putnam--Gold)} (trial);
\end{tikzpicture}
\caption[The identification hierarchy as a chain of separations]{The identification hierarchy, read as a
chain of strict separations refined by a lattice of single-constraint relaxations. Each spine node
($\FIN$, $\Ex$, $\BC$) is a success criterion; each spine edge is a separation theorem, a class one
criterion learns and the other cannot, not a more liberal definition. The bottom rung
$\FIN\subsetneq\Ex$ is verified in the kernel (\leanname{ex_strictly_stronger_finite}), drawn solid;
the witness is the class of all finite languages, which is $\Ex$-identifiable but not
$\FIN$-identifiable. The upper rung $\Ex\subsetneq\BC$ is the Case--Smith separation, where $\Ex$
demands the index converge (syntactic) and $\BC$ demands only the language converge (semantic); it is
established in the literature and not yet in the kernel, so it is drawn dashed and cited. The side nodes
are single-constraint relaxations of $\Ex$: anomalous $\Ex^*$ removes the zero-error constraint and so
extends $\Ex$; monotonic $\mathrm{Mon}\Ex$ adds a no-retraction constraint and so restricts $\Ex$;
vacillatory $\mathrm{Vex}$ allows finitely many final hypotheses and sits strictly between $\Ex$ and
$\BC$; trial-and-error learning pins the limit-computable predicates to the $\Delta^0_2$ level. Every
branch and the $\Delta^0_2$ link are literature, not kernel theorems, hence dashed.}
\label{fig:identification-hierarchy}
\end{figure}

\begin{traversal}
\textbf{Path:} $\nde{fin\_learning} \xrightarrow{\rel{strictly\_stronger}} \nde{ex\_learning} \xrightarrow{\rel{strictly\_stronger}} \nde{bc\_learning}$.\\[3pt]
Each arrow is witnessed by an explicit separation. The graph encodes these as \rel{strictly\_stronger} edges with the separation witness as metadata. The hierarchy is not a sequence of increasingly liberal definitions; it is a chain of \emph{theorems}, each proved by constructing a class that one criterion can learn and the other cannot.
\end{traversal}

% ================================================================
% SECTION 4: Relaxations and Their Lattice
% ================================================================

\section{Relaxations of Identification}
\label{sec:relaxations}

The $\FIN$--$\Ex$--$\BC$ chain is the spine of the hierarchy. Several natural relaxations branch off from $\Ex$, each modifying a different aspect of the convergence requirement.

\subsection{Anomalous Learning}
\label{sec:anomalous-gold}

\begin{defn}[Anomalous Learning ($\Ex^*$)]
\label{def:anomalous}
A learner \emph{$\Ex^*$-identifies} a language $L$ if it $\Ex$-converges to an index $e$ such that $W_e$ differs from $L$ on at most finitely many strings. That is, the final hypothesis may contain finitely many errors: finitely many strings incorrectly included or excluded.\formal{FLT_Proofs.Criterion.Gold.AnomalousLearnable}
\end{defn}

Anomalous learning relaxes $\Ex$ by removing the requirement of exact correctness, replacing it with correctness up to a finite set. In graph terms, this is a \rel{restricts} edge from $\nde{anomalous\_learning}$ to $\nde{ex\_learning}$: the constraint (zero anomalies) is removed, and no new grammatical structure is introduced. The class of $\Ex^*$-identifiable families strictly contains the class of $\Ex$-identifiable families~\cite{CaseSmith1983}.

\subsection{Monotonic Learning}
\label{sec:monotonic}

\begin{defn}[Monotonic Learning ($\mathrm{Mon}\Ex$)]
\label{def:monotonic}
A learner \emph{monotonically} $\Ex$-identifies $L$ if it $\Ex$-identifies $L$ and, whenever it changes its hypothesis from $e$ to $e'$, the new hypothesis is at least as inclusive on the data seen so far: $W_e \cap \{t_0, \ldots, t_n\} \subseteq W_{e'} \cap \{t_0, \ldots, t_n\}$. Once a datum is correctly classified, that classification is never retracted.\formal{FLT_Proofs.Criterion.Gold.MonotonicLearnable}
\end{defn}

Monotonic learning \emph{strengthens} $\Ex$ by adding a constraint. The class of $\mathrm{Mon}\Ex$-identifiable families is a strict subset of $\Ex$-identifiable families. In the graph, $\nde{monotonic\_learning} \xrightarrow{\rel{restricts}} \nde{ex\_learning}$.

\subsection{Vacillatory Learning}
\label{sec:vacillatory}

\begin{defn}[Vacillatory Learning ($\mathrm{Vex}$)]
\label{def:vacillatory}
A learner \emph{vacillatorily} identifies $L$ if it eventually oscillates among finitely many indices $e_1, \ldots, e_k$, all satisfying $W_{e_i} = L$. The learner never settles on a single index, but all its eventual outputs are extensionally correct and drawn from a finite set.\formal{FLT_Proofs.Criterion.Gold.VacillatoryLearnable}
\end{defn}

Vacillatory learning sits between $\Ex$ (which requires convergence to a single index) and $\BC$ (which allows infinitely many correct indices). It is strictly more powerful than $\Ex$ and strictly less powerful than $\BC$.

\subsection{Trial and Error}
\label{sec:trial-error}

\begin{defn}[Trial-and-Error Learning]
\label{def:trial-error}
A \emph{trial-and-error} predicate for a set $A \subseteq \N$ is a computable function $f : \N \to \{0, 1\}$ such that $\lim_{s \to \infty} f_s(n)$ exists for all $n$ and equals the characteristic function of $A$. The learner may ``retract'' previous outputs: it learns by making mistakes and correcting them.\formal{FLT_Proofs.Criterion.Gold.TrialAndErrorLearnable}
\end{defn}

Trial-and-error learning connects identification in the limit to the arithmetical hierarchy. A set $A$ is trial-and-error learnable if and only if $A \in \Delta^0_2$, the class of sets that are both $\Sigma^0_2$ and $\Pi^0_2$. This is the precise recursion-theoretic characterization, due to Putnam~\cite{Gold1965} and Gold: the limit-computable functions are exactly the $\Delta^0_2$ functions.\formal{FLT_Proofs.Computation.Delta02Class} This connection places Gold-style learning firmly within the landscape of classical recursion theory.

% ================================================================
% SECTION 5: Mind-Change Complexity and Ordinals
% ================================================================

\section{Mind-Change Complexity}
\label{sec:mind-change}

Gold's impossibility theorem tells us that certain classes cannot be identified at all. For classes that \emph{can} be identified, a natural quantitative question arises: how many times must the learner change its mind before converging?

\begin{defn}[Mind Change]
\label{def:mind-change}
A \emph{mind change} occurs at time $n$ if the learner's output satisfies $M(t_0, \ldots, t_n) \neq M(t_0, \ldots, t_{n-1})$. The \emph{mind-change count} of $M$ on text $t$ is the number of times $M$ changes its hypothesis.\formal{FLT_Proofs.Complexity.MindChange.MindChangeCount}
\end{defn}

For finite mind-change bounds, the theory is straightforward: we say that $M$ identifies $\mathcal{L}$ with at most $k$ mind changes if, on every text for every $L \in \mathcal{L}$, the mind-change count is at most $k$. The class of all finite languages is identifiable with 0 mind changes after the first hypothesis (the learner in \Cref{ex:finite-languages} changes its mind every time a new element appears, but a more careful learner can be designed with bounded mind changes for restricted subclasses).

The surprise (genuinely unexpected for readers coming from PAC theory) is that integer-valued bounds are \emph{not sufficient} to characterize the full landscape.

\begin{thm}[Freivalds--Smith {\cite{FreivaldsSmith1993}}]
\label{thm:mind-change-ordinal}
The mind-change complexity of $\Ex$-identification is naturally measured by \emph{countable ordinals}. Specifically:
\begin{enumerate}[label=(\roman*),leftmargin=*]
\item For every countable ordinal $\alpha$, one can define what it means for a learner to identify a class with mind-change bound $\alpha$.
\item There exist classes identifiable with mind-change bound $\omega$ (the first infinite ordinal) that cannot be identified with any finite mind-change bound.
\item More generally, for every ordinal $\alpha < \omega_1$, there exist classes identifiable with mind-change bound $\alpha$ but not with any bound $\beta < \alpha$.
\end{enumerate}\formal{FLT_Proofs.Complexity.MindChange.MindChangeOrdinal}\formal{FLT_Proofs.Theorem.Gold.mind_change_characterization}
\end{thm}

The idea behind ordinal mind-change bounds is as follows. A learner with mind-change bound $\omega$ begins with an ordinal counter set to $\omega$. At each mind change, the counter must decrease, but it may decrease to any smaller ordinal. Since there is no infinite descending sequence of ordinals (ordinals are well-ordered), the learner must eventually stop changing its mind. The counter $\omega$ allows finitely many mind changes, but the \emph{number} of mind changes need not be bounded in advance by any fixed integer; it may depend on the input.

This is fundamentally different from an integer bound. With a bound of $k \in \N$, the learner may change its mind at most $k$ times on \emph{every} input. With a bound of $\omega$, the learner may change its mind $k$ times for input-dependent $k$, with no uniform finite upper bound. The ordinal hierarchy continues: $\omega \cdot 2$ allows the learner to ``reset'' its finite counter once; $\omega^2$ allows nested levels of resetting; and so on up through the constructive ordinals.

\begin{remark}[Ordinals in learning theory]
\label{rmk:ordinals-surprise}
The appearance of transfinite ordinals in learning theory is a genuine structural surprise. PAC learning theory uses real-valued parameters ($\varepsilon, \delta, m(\varepsilon, \delta)$). Online learning uses integer-valued dimensions ($\Ldim$). Gold-style learning requires ordinal-valued complexity measures. Each proof technique brings its own number system. The full development of ordinal mind-change complexity is deferred to \Cref{ch:ordinals}, where it interacts with the constructive ordinal notation systems of Kleene.
\end{remark}

\begin{traversal}
\textbf{Path:} $\nde{mind\_change\_characterization} \xrightarrow{\rel{measures}} \nde{ex\_learning}$.\\[3pt]
The mind-change ordinal is a complexity measure on $\Ex$-identifiable classes, analogous to the VC dimension for PAC-learnable classes. But where VC dimension is a single integer, mind-change complexity is an ordinal, potentially transfinite. This is a \rel{measures} edge of a qualitatively different kind than $\nde{vc\_dimension} \xrightarrow{\rel{measures}} \nde{concept\_class}$.
\end{traversal}

% ================================================================
% SECTION 6: The Three-Paradigm Separation
% ================================================================

\subsection{In-context learning as identification in the limit}
\label{sec:icl-gold}

Gold's framework (infinite horizon, a stream of data, a sequence of conjectures) describes a learning mode that modern transformers exhibit directly: in-context learning.  A transformer fed a growing prompt of labelled examples produces, at each prefix length, a prediction rule: its conjecture given what it has seen so far.  The conjecture map ``finite data so far $\mapsto$ predictor'' is exactly the type of a Gold learner.

\begin{prop}[A transformer in context is a Gold learner]
\label{prop:transformer-gold}
A parametric attention learner $T$, read as the map from a finite prefix $((x_1,y_1),\dots,(x_n,y_n))$ to the predictor it computes in context, is a Gold learner: a function from finite labelled sequences to concepts.\formal{FLT_Proofs.Learner.Core.GoldLearner} Its behaviour on a target and data stream is therefore classified by the criteria of this chapter ($\Ex$, $\BC$, vacillatory), and its mind-change complexity is the ordinal-valued measure of \cref{thm:mind-change-ordinal} applied to $T$.\formal{FLT_Proofs.Complexity.MindChange.MindChangeOrdinal}
\end{prop}

\begin{proof}
The Gold-learner type is a conjecture map from $\mathrm{List}(X \times Y)$ to concepts, and the in-context predictor of $T$ at prefix $s$ is such a map: feed $s$ as context and return the function $x \mapsto T(s)(x)$.  The mind-change count and ordinal of \cref{def:mind-change,thm:mind-change-ordinal} are defined for every Gold learner, so they apply to $T$ verbatim.
\end{proof}

The proposition turns a practical question (does in-context learning ``converge''?) into the precise question of which identification criterion $T$ meets.  Syntactic convergence, where the in-context predictor stops changing, is $\Ex$; extensional convergence, where the predictor keeps changing but computes a fixed function, is $\BC$; oscillation among finitely many correct predictors is vacillatory.  Gold's impossibility theorem then bites: no transformer $\Ex$-identifies, from positive data alone, a class containing all finite concepts and one infinite concept (\cref{thm:gold}).  Whether realistic in-context learning $\Ex$-converges, $\BC$-converges, or vacillates, and with what mind-change ordinal, is open (\cref{op:icl-identification}).

\section{Three Paradigms: Online Inside PAC, Gold Apart}
\label{sec:three-paradigm}

We now have three learning paradigms: PAC (\Cref{ch:pac}), online (\Cref{ch:online}), and Gold. Each defines ``learnable'' differently. The question is: what is the relationship?

The answer is a partial order, not an antichain. Online learnability is \emph{strictly contained} in PAC learnability: every online-learnable class is PAC-learnable, and the converse fails, so the two are comparable. Gold-style identification, by contrast, is incomparable to both: it identifies classes of unbounded VC and Littlestone dimension that no statistical learner can handle, while it fails on classes with no convergence signal that PAC and online learners handle easily. The deep structural fact is that statistical capacity (PAC and online) and convergence in the limit (Gold) are orthogonal axes, and that adaptivity, the only thing separating online from PAC, buys nothing here: a finite Littlestone dimension forces a finite VC dimension, so online sits properly inside PAC.

\begin{separation}
\textbf{Theorem} (PAC, online, and Gold).\formal{FLT_Proofs.Theorem.Separation.online_pac_gold_separation} \textit{Online learnability is strictly contained in PAC learnability, so the two are comparable; Gold-style identification is incomparable to both. The machine-checked statement establishes the threshold separation (PAC but not online) and separation (a) below (the finite languages, $\Ex$ but not PAC); the remaining separations (b)--(d) are classical:}

\begin{enumerate}[label=(\alph*),leftmargin=*,itemsep=4pt]
\item \textbf{$\Ex$-learnable but not PAC-learnable.} The class $\mathcal{L}_{\mathrm{fin}}$ of all finite subsets of $\N$ is $\Ex$-identifiable from text (\Cref{ex:finite-languages}). But the associated concept class has infinite VC dimension (any finite set of points can be shattered), so it is not PAC-learnable.

The witness exploits the fundamental difference in how the two paradigms treat ``all finite subsets.'' $\Ex$-identification succeeds because on any \emph{fixed} finite set, the learner eventually sees all elements. PAC fails because the learner must handle \emph{all} finite sets simultaneously with bounded sample complexity, and the VC dimension is infinite.

\item \textbf{PAC-learnable but not $\Ex$-identifiable.} The class of threshold functions $\mathcal{C} = \{x \mapsto \ind[x \geq \theta] : \theta \in \R\}$ over $\R$ is PAC-learnable ($\VC = 1$). But it is not $\Ex$-identifiable from text, because a text for a threshold language $\{\theta, \theta+1, \theta+2, \ldots\}$ reveals the threshold only in the limit, and the class of all such languages together with all finite languages triggers Gold's impossibility theorem.

\item \textbf{Online-learnable but not $\Ex$-identifiable.} The class of singletons $\mathcal{C} = \{\{x\} : x \in X\}$ has Littlestone dimension~1 (online-learnable with at most 1 mistake). But the class of all singletons together with the empty set is not $\Ex$-identifiable from text for reasons analogous to Gold's theorem: the learner cannot distinguish ``no more positive examples will come'' from ``the next positive example has not yet arrived.''

\item \textbf{$\Ex$-identifiable but not online-learnable.} Consider any class $\mathcal{L}$ of recursive languages that is $\Ex$-identifiable and has infinite Littlestone dimension. Such classes exist: the class of all pattern languages (Angluin, 1980) is $\Ex$-identifiable from text but has infinite Littlestone dimension over suitable domains.
\end{enumerate}
\end{separation}

\begin{figure}[ht]
\centering
\begin{tikzpicture}[
    >={Stealth[length=2mm]},
    territory/.style={draw, thick, rounded corners=8pt, align=center, font=\small},
    wit/.style={font=\scriptsize, text=edgeorange, fill=white, inner sep=1.5pt, align=center},
    hdl/.style={font=\tiny, text=black!65, align=center},
    imp/.style={-{Stealth[length=2.2mm]}, very thick, defblue!60!black},
    sepv/.style={-{Stealth[length=2mm]}, thick, thmred},
    sepl/.style={-{Stealth[length=2mm]}, thick, thmred!85, densely dashed}
]
% --- PAC outer region; Online nested inside it; Gold a separate region ---
\node[territory, fill=defblue!8, minimum width=5.2cm, minimum height=4.0cm] (pac) at (0,0) {};
\node[font=\small\bfseries, defblue!75!black, anchor=north west, align=left]
   at ([xshift=5pt,yshift=-4pt]pac.north west) {PAC\\[-1pt]{\scriptsize\mdseries $\varepsilon,\delta$; VC dim}};
\node[territory, fill=onlineblue!14, minimum width=3.0cm, minimum height=1.7cm] (online) at (0.45,-0.7)
   {Online\\[-1pt]{\scriptsize\mdseries Littlestone dim}};
\node[territory, fill=goldcolor!12, minimum width=3.2cm, minimum height=4.0cm] (gold) at (7.6,0)
   {Gold / $\Ex$\\[-1pt]{\scriptsize\mdseries identification}\\{\scriptsize\mdseries in the limit}\\[3pt]{\scriptsize\mdseries ordinal /}\\{\scriptsize\mdseries mind-change}};
% --- E1: Online => PAC (verified containment), strict ---
\draw[imp] (0.45,0.15) -- (0.45,1.5);
\node[hdl, anchor=south, text=defblue!65!black] at (0.45,1.52) {\leanname{online_imp_pac}};
\node[hdl, anchor=north] at (0.45,-1.7)
   {strict: thresholds $\VC{=}1$, $\Ldim{=}\infty$\\{\tiny\leanname{pac_not_implies_online}}};
% --- E2: Gold =/=> PAC (verified separation, solid) ---
\draw[sepv] (6.0,1.3) -- (2.6,1.3);
\node[wit, fill=white] at (4.3,1.3) {finite sets: $\Ex$ yes, PAC no\\{\tiny\leanname{ex_not_implies_pac}}};
% --- E3,E4,E5: literature separations (dashed) ---
\draw[sepl] (2.6,0.45) -- (6.0,0.45);
\node[wit, fill=white] at (4.3,0.45) {thresholds: PAC yes, $\Ex$ no\\{\tiny\cite{Gold1967}}};
\draw[sepl] (1.95,-0.55) -- (6.0,-0.55);
\node[wit, fill=white] at (4.3,-0.55) {singletons$+\emptyset$: Online yes, $\Ex$ no\\{\tiny\cite{Gold1967}}};
\draw[sepl] (6.0,-1.4) -- (1.95,-1.4);
\node[wit, fill=white] at (4.3,-1.4) {patterns: $\Ex$ yes, Online no\\{\tiny\cite{Angluin1980}}};
\end{tikzpicture}
\caption[The three paradigms: online strictly inside PAC, Gold incomparable to both]{The relationship
among the three paradigms, corrected to the verified structure: it is a partial order, not an antichain.
Online learnability is strictly contained in PAC learnability: every online-learnable class is PAC-learnable
(\leanname{online_imp_pac}, solid), and the containment is strict because the threshold class over $\R$ is
PAC-learnable but has infinite Littlestone dimension, hence is not online-learnable
(\leanname{pac_not_implies_online}); so the Online region sits inside the PAC region. There is no class that
is online-learnable but not PAC-learnable, and none with finite Littlestone but infinite VC dimension, since
$\VC\le\Ldim$. Gold-style identification is incomparable to both: it identifies the finite languages, which
have infinite VC dimension and so escape every statistical learner (\leanname{ex_not_implies_pac}, solid),
while it cannot identify classes with no convergence signal that PAC and online learners handle (thresholds,
singletons with the empty set; dashed, \cite{Gold1967}), and pattern languages give the reverse (dashed,
\cite{Angluin1980}). Solid arrows are kernel theorems; dashed arrows are literature.}
\label{fig:three-paradigm}
\end{figure}

\begin{remark}[What the separations mean]
\label{rmk:separation-meaning}
The three-paradigm separation is not a deficiency of the definitions. It reflects a genuine mathematical fact: different notions of ``learning from data'' capture different aspects of the learning problem, and no single notion is universal. PAC learning measures statistical generalization under distributional assumptions. Online learning measures worst-case sequential prediction. Gold-style identification measures eventual convergence without any accuracy or efficiency guarantee. These are three different mathematical questions, and they have three different answers.
\end{remark}

% ================================================================
% SECTION 7: Chapter Summary
% ================================================================

\section{What This Chapter Established}
\label{sec:ch7-summary}

This chapter introduced Gold-style identification in the limit, the oldest paradigm in formal learning theory, and established five structural facts:

\begin{enumerate}[leftmargin=*,itemsep=3pt]
\item \textbf{Gold's impossibility theorem} (\Cref{thm:gold}): any class containing all finite languages and at least one infinite language is not $\Ex$-identifiable from text. The proof is by diagonalization, the first impossibility result in learning theory and one whose proof technique (adversarial stream construction) has no analogue in PAC or online theory.

\item \textbf{The identification hierarchy} (\Cref{sec:hierarchy}): the chain $\FIN \subsetneq \Ex \subsetneq \BC$ is strict, with the Case--Smith separation (\Cref{thm:bc-strictly-stronger}) providing the witness for $\BC > \Ex$. The hierarchy is not a sequence of definitions but a chain of theorems.

\item \textbf{Relaxations form a lattice} (\Cref{sec:relaxations}): anomalous, monotonic, and vacillatory learning branch off from $\Ex$, each modifying a different aspect of the convergence requirement. Trial-and-error learning connects to $\Delta^0_2$ in the arithmetical hierarchy.

\item \textbf{Mind-change complexity is ordinal-valued} (\Cref{sec:mind-change}): the Freivalds--Smith characterization shows that the right complexity measure for $\Ex$-identification uses transfinite ordinals, not integers. The full treatment is in \Cref{ch:ordinals}.

\item \textbf{The three paradigms form a partial order, not an antichain} (\Cref{sec:three-paradigm}): online learnability is strictly contained in PAC learnability (every online-learnable class is PAC-learnable, and the converse fails), while Gold-style identification is incomparable to both. Statistical capacity and convergence in the limit are orthogonal axes.
\end{enumerate}

The recursion-theoretic character of this chapter (diagonalization proofs, connections to the arithmetical hierarchy, ordinal-valued measures) contrasts sharply with the probabilistic character of \Cref{ch:pac} and the combinatorial character of \Cref{ch:online}. Each paradigm brings its own mathematical world. The separations of \Cref{sec:three-paradigm} are consequences of this fact: the paradigms use different proof techniques because they formalize genuinely different questions.


% ================================================================
% EXERCISES
% ================================================================

\subsection*{Exercises}

\begin{enumerate}[leftmargin=*,itemsep=6pt]

\item \textbf{The power of informant over text.}
  Gold showed that the class $\mathcal{L}$ of all recursive languages
  is $\Ex$-identifiable from \emph{informant} (both positive and
  negative examples) but not from text (positive examples only).
  \begin{enumerate}[label=(\alph*)]
  \item Prove the informant direction: construct an $\Ex$-learner $M$
    that identifies any recursive language $L$ from informant.
    \emph{Hint:}  Dovetail over all programs $\varphi_0, \varphi_1,
    \ldots$, and at time $t$ hypothesize the smallest index~$e$
    consistent with all labeled data seen so far.  Use the fact that
    the informant eventually presents every string with its correct
    label, so incorrect indices are eventually refuted.
  \item Explain precisely why this learner fails from text.  Identify
    the step in the argument that relies on negative examples, and
    show that no text-based substitute exists.  (\emph{Hint:} The
    learner can refute ``$\varphi_e$ includes $x$'' from an
    informant presentation of $(x, 0)$, but a text for $L$ never
    explicitly says ``$x \notin L$''; it merely fails to present
    $x$, which is indistinguishable from delay.)
  \end{enumerate}

\item \textbf{$\Ex$-identification of co-finite languages.}
  Let $\mathcal{L}_{\mathrm{cof}}$ be the class of all co-finite
  languages over $\N$: $L \in \mathcal{L}_{\mathrm{cof}}$ iff
  $\N \setminus L$ is finite.
  \begin{enumerate}[label=(\alph*)]
  \item Prove that $\mathcal{L}_{\mathrm{cof}}$ is
    $\Ex$-identifiable from text.  (\emph{Hint:} Every co-finite
    language $L$ contains all but finitely many natural numbers.
    A text for $L$ eventually presents every element of $L$, so after
    time $n_0$, every natural number $\leq n_0$ has either appeared
    in the text (and is in $L$) or has not appeared (and may or
    may not be in $L$).  Design a learner that waits long enough
    to conclude that unseen small numbers are \emph{not} in $L$.)
  \item Prove that $\mathcal{L}_{\mathrm{fin}} \cup
    \mathcal{L}_{\mathrm{cof}}$ (all finite and all co-finite
    languages together) is \emph{not} $\Ex$-identifiable from text.
    (\emph{Hint:} Apply Gold's impossibility theorem.  Verify
    that this class contains all finite languages and at least one
    infinite language.)
  \item The result of~(b) is sharper than Gold's theorem applied
    to ``all finite + one infinite'': the single infinite language
    is itself very structured (co-finite).  Explain why the structure
    of the infinite language does not help: what makes the
    diagonalization work is not the complexity of $L_\infty$ but
    the learner's inability to distinguish ``finite set, done
    growing'' from ``co-finite set, still growing.''
  \end{enumerate}

\end{enumerate}

\section{Open Problems}
\label{sec:ch7-open}

Each problem is anchored to a result above and tagged: a \emph{known unknown} (\textsc{ku}) is an articulated open question; an \emph{unknown known} (\textsc{uk}) is a result in the literature or the verified development not yet absorbed into a clean statement.  Verified handles are named where they exist.

\begin{openprob}[The identification criterion of in-context learning, \textsc{uk}]
\label{op:icl-identification}
\Cref{prop:transformer-gold} casts a transformer's in-context predictor as a Gold learner.  Determine which criterion it meets on natural concept classes: does a fixed transformer $\Ex$-converge (the in-context rule stabilizes), $\BC$-converge (it stabilizes only extensionally), or vacillate, and what is its mind-change ordinal (\leanname{FLT_Proofs.Complexity.MindChange.MindChangeOrdinal}) as a function of context length and model size?
\end{openprob}

\begin{openprob}[Formalizing the $\BC > \Ex$ separation, \textsc{uk}]
\label{op:bc-ex-formal}
Gold's impossibility theorem is verified (\leanname{FLT_Proofs.Theorem.Gold.gold_theorem}) and the criterion hierarchy is typed, but the Case--Smith strict separation (\cref{thm:bc-strictly-stronger}) is not.  Formalize the construction of a class that is $\BC$-identifiable but not $\Ex$-identifiable, on top of \leanname{FLT_Proofs.Criterion.Gold.BCLearnable} and \leanname{FLT_Proofs.Criterion.Gold.EXLearnable}.
\end{openprob}

\begin{openprob}[The Freivalds--Smith ordinal hierarchy, \textsc{uk}]
\label{op:fs-ordinal}
\Cref{thm:mind-change-ordinal} is stated but its strictness is cited.  Formalize the transfinite mind-change hierarchy on top of the typed ordinal measure (\leanname{FLT_Proofs.Complexity.MindChange.MindChangeOrdinal}): prove that for each constructive ordinal $\alpha$ the $\alpha$-bounded classes strictly contain the $\beta$-bounded classes for $\beta < \alpha$, connecting to the constructive ordinal notations of \Cref{ch:ordinals}.
\end{openprob}

\begin{openprob}[Angluin's telltale characterization, \textsc{ku}]
\label{op:telltale}
Gold's theorem is an impossibility result; the matching characterization is Angluin's: a class is $\Ex$-identifiable from text if and only if every language in it has a finite ``telltale'' subset that no proper sublanguage in the class contains~\cite{Angluin1980}.  Formalize this characterization on top of \leanname{FLT_Proofs.Criterion.Gold.EXLearnable}, giving the exact condition that separates the identifiable classes from those Gold's theorem rules out.
\end{openprob}

\begin{openprob}[The vacillatory hierarchy, \textsc{ku}]
\label{op:vex-hierarchy}
Vacillatory identification (\cref{def:vacillatory}) allows oscillation among at most $k$ correct indices.  Determine whether the $\mathrm{Vex}_k$ hierarchy is strict in $k$ and where $\mathrm{Vex} = \bigcup_k \mathrm{Vex}_k$ sits strictly between $\Ex$ and $\BC$, formalizing on top of \leanname{FLT_Proofs.Criterion.Gold.VacillatoryLearnable}.
\end{openprob}

\begin{openprob}[A meta-measure across the three paradigms, \textsc{ku}]
\label{op:meta-measure}
The three-paradigm structure (\cref{sec:three-paradigm}, \leanname{FLT_Proofs.Theorem.Separation.online_pac_gold_separation}) places online learnability strictly inside PAC and Gold-style identification incomparable to both.  Determine whether a single parameterized complexity measure specializes to the VC dimension, the Littlestone dimension, and the mind-change ordinal under three settings of a data-model parameter, or prove that the three invariants are provably non-unifiable.
\end{openprob}

\begin{openprob}[Anomalous in-context learning, \textsc{uk}]
\label{op:anomalous-icl}
Anomalous identification (\cref{def:anomalous}, \leanname{FLT_Proofs.Criterion.Gold.AnomalousLearnable}) permits finitely many errors in the final hypothesis.  Determine whether relaxing exactness to bounded anomaly enlarges what a transformer can identify in context, and relate the anomaly count to the approximation error of the attention learner on the tail of the data stream.
\end{openprob}

\begin{openprob}[Identification from informant versus text for attention, \textsc{ku}]
\label{op:informant-attention}
Exercise~1 shows informant is strictly more powerful than text.  Quantify the gap for a transformer: how much does access to negative examples (an informant-style prompt) reduce the mind-change ordinal of in-context identification relative to a positive-only (text-style) prompt, and is the separation as sharp for bounded-context transformers as it is in the recursive setting?
\end{openprob}

\begin{openprob}[Monotonic in-context learning, \textsc{uk}]
\label{op:monotonic-icl}
Monotonic identification (\cref{def:monotonic}, \leanname{FLT_Proofs.Criterion.Gold.MonotonicLearnable}) forbids retracting a correct classification.  Determine whether attention learners can be constrained to be monotonic without loss of identification power on natural classes, and whether monotonicity corresponds to a stability property of the routing under prefix extension.
\end{openprob}

\begin{openprob}[Effective three-paradigm witnesses, \textsc{uk}]
\label{op:effective-three-paradigm}
The three-paradigm separation produces its witnesses existentially.  Give a single parameterized family of concept classes that realizes each of the separations of \cref{sec:three-paradigm} (online strictly inside PAC, Gold incomparable to both) with explicit witnesses and explicit obstructions, turning the separation lattice into one constructive theorem rather than several independent constructions.
\end{openprob}


\subsection*{Counterfactual exercises: generalizing Figure~\ref{fig:identification-hierarchy}}

Each exercise perturbs the hierarchy. Show, in the figure's grammar (a strict-power partial order on identification criteria whose every edge is a separation theorem, refined by a relaxation lattice of single-constraint edits to the success notion, solid where a kernel theorem certifies the separation and dashed where it is borrowed from the literature or open), that the generalized picture subsumes the concrete one as an \emph{instance} (an edge whose separation is determined), a \emph{boundary} (where one constraint is toggled and the class strictly grows or shrinks), or a \emph{frontier} (a dashed, open or literature-only edge).

\begin{enumerate}[leftmargin=*,itemsep=3pt]
\item \textbf{The one verified rung.} Read the bottom edge $\FIN \to \Ex$ not as a containment of definitions but as a theorem with a witness class. Show this is the figure's only solid instance: the kernel proves $\FIN$-learnability implies $\Ex$-learnability by dropping the no-mind-change conjunct (\leanname{ex_strictly_stronger_finite}), and the class of all finite languages is the separating witness (\Cref{ex:finite-languages}), so this rung alone is certified rather than cited.

\item \textbf{Syntactic versus semantic convergence is the whole gap.} Replace the labels $\Ex$ and $\BC$ by their quantifier skeletons. Show this is the boundary that names the upper rung: $\Ex$ requires the index to satisfy $h = c$ (\leanname{EXLearnable}) while $\BC$ requires only $\forall x,\ h\,x = c\,x$ (\leanname{BCLearnable}), one clause toggled from equality of programs to equality of extensions, and that single toggle is exactly what the Case--Smith class exploits.

\item \textbf{Why the upper rung is dashed.} Ask the kernel for the $\Ex \subsetneq \BC$ separation. Show this is the figure's principal frontier: the criteria are typed (\leanname{EXLearnable}, \leanname{BCLearnable}) but no theorem separates them, so the edge is drawn dashed and cited (\cite{CaseSmith1983}, \cref{op:bc-ex-formal}); contrast this with the solid rung of Exercise~1 to see that line style is evidence grade, not decoration.

\item \textbf{Anomalous learning extends, not restricts.} Toggle the exactness clause of $\Ex$ from $h = c$ to ``$h$ errs on at most $a$ strings''. Show this is a boundary on the larger side: removing a constraint can only enlarge the learnable class, so $\Ex^*$ \rel{extends} $\Ex$ (\leanname{AnomalousLearnable}), the separation is strict but literature-only (\cite{CaseSmith1983}, \cref{op:anomalous-icl}), hence the $\Ex^* \to \Ex$ edge is dashed and oriented as an extension.

\item \textbf{Monotonic learning restricts.} Conjoin to $\Ex$ the no-retraction clause $h\,\mathrm{data}\,x = \mathrm{true} \Rightarrow h\,(\mathrm{data}\,{+}{+}\,[xy])\,x = \mathrm{true}$. Show this is the boundary on the smaller side: adding a constraint can only shrink the class, so $\mathrm{Mon}\Ex$ \rel{restricts} $\Ex$ (\leanname{MonotonicLearnable}); the strictness is literature (\cite{OshersonStobWeinstein1986}, \cref{op:monotonic-icl}), so the edge is dashed and oriented as a restriction, the mirror image of Exercise~4.

\item \textbf{Vacillation is two toggles at once.} Replace the single final index by membership in a finite set of correct indices. Show this is the instance that sits the branch between the spine nodes: relaxing $\Ex$'s one index to finitely many extends $\Ex$, while restricting $\BC$'s arbitrarily many indices to finitely many is the dual restriction of $\BC$, so $\mathrm{Vex}$ (\leanname{VacillatoryLearnable}) lies strictly between them, with both legs dashed and cited (\cite{CaseSmith1983}, \cref{op:vex-hierarchy}).

\item \textbf{The arithmetical anchor.} Read the $\Delta^0_2$ tag on the trial-and-error node as a characterization, not a containment. Show this is a frontier that re-types the value: trial-and-error learnability is pointwise convergence (\leanname{TrialAndErrorLearnable}) and equals membership in $\Delta^0_2$ (\leanname{Delta02Class}) by the Putnam--Gold theorem, but the kernel carries the class as a definition and not the iff, so the characterization edge is dashed and cited (\cite{Putnam1965}).

\item \textbf{Gold's theorem is the pressure behind the chain.} Ask why the hierarchy branches at $\Ex$ at all. Show this is the instance that motivates every edge: Gold's impossibility (\leanname{gold_theorem}) proves $\neg\,\Ex$-identifiability for any class with all finite languages and one infinite language, so strengthening to $\FIN$ and weakening to $\BC$ are the two escape routes the figure draws, and the chain exists because $\Ex$ is provably bounded.

\item \textbf{Mind changes turn the spine into an ordinal.} Refine the $\FIN$-to-$\Ex$ gap by counting retractions rather than forbidding them. Show this is a frontier that grades the bottom rung: $\FIN$ is zero mind changes and $\Ex$ is finitely many, and the kernel types the intermediate measure (\leanname{MindChangeOrdinal}), but the strictness of the transfinite hierarchy between them is cited, not proved (\cite{FreivaldsSmith1993}, \cref{op:fs-ordinal}), so the graded edge is dashed.

\item \textbf{Is the lattice one graded criterion?} Ask whether all seven nodes are evaluations of a single success criterion parameterized by four toggles: error budget, retraction flag, index cardinality, and convergence mode (syntactic or semantic). Show this is the figure's most general open frontier: it would present the hierarchy not as a fixed chain plus branches but as one criterion space with the four constraints as a basis, yet the kernel formalizes the seven criteria separately with no graded predicate unifying them, so the meta-edge is dashed over an otherwise typed lattice.
\end{enumerate}

\subsection*{Counterfactual exercises: generalizing Figure~\ref{fig:three-paradigm}}

Each exercise perturbs the corrected picture. Show, in the figure's grammar (a strict-power partial order on learning paradigms whose every edge is either an implication with a strict witness or a separation with a witness, solid where a kernel theorem certifies it and dashed where it is borrowed from the literature), that the generalized picture subsumes the concrete one as an \emph{instance} (an edge whose direction and evidence grade are determined), a \emph{boundary} (where one constraint is toggled and the learnable class strictly grows or shrinks), or a \emph{frontier} (a dashed, open or literature-only edge).

\begin{enumerate}[leftmargin=*,itemsep=3pt]
\item \textbf{The picture is an order, not an antichain.} Read the three paradigms not as three incomparable territories but as nodes of a partial order under ``strictly more powerful than''. Show this is the figure's organizing instance: online learnability implies PAC learnability (\leanname{online_imp_pac}) and the converse fails (\leanname{pac_not_implies_online}), so the Online and PAC nodes are \emph{comparable} and the whole picture is a containment plus an incomparable Gold, never a symmetric triangle of six arrows.

\item \textbf{The impossible witness.} Try to draw the edge ``some class is online-learnable but not PAC-learnable'', certified by ``finite Littlestone dimension, infinite VC dimension''. Show this is the frontier that does not exist: $\VC \le \Ldim$ (\leanname{BranchWiseLittlestoneDim_ge_VCDim}) makes the certificate ``$\Ldim<\infty$ and $\VC=\infty$'' an empty predicate, and \leanname{online_imp_pac} makes ``online but not PAC'' an empty class, so this edge is removed rather than dashed. A separation drawn with an impossible witness is not a weak edge; it is no edge.

\item \textbf{Strictness is the threshold class.} Read the strictness of $\text{Online}\subsetneq\mathrm{PAC}$ as a theorem with a witness, not a gap. Show this is the boundary on the larger side: the threshold class over $\N$ is PAC-learnable yet has $\Ldim=\infty$, hence is not online-learnable (\leanname{pac_not_implies_online}), so toggling from ``i.i.d.\ sample'' to ``adversarial order'' strictly shrinks the class, and this single witness is what makes the Online region sit \emph{properly} inside PAC.

\item \textbf{One solid arrow certifies two separations.} Ask the kernel only for $\Ex\not\Rightarrow\mathrm{PAC}$. Show this is the instance that does double duty: the finite-languages witness is $\Ex$-identifiable but has $\VC=\infty$ (\leanname{ex_not_implies_pac}), and because Online lies inside PAC, the same class is a fortiori not online-learnable, so the lone solid arrow $\Ex\not\Rightarrow\mathrm{PAC}$ simultaneously certifies $\Ex\not\Rightarrow\text{Online}$ and the cluster $\{\mathrm{PAC},\text{Online}\}$ needs no separate Gold-escapes-Online theorem.

\item \textbf{What the three-way theorem actually says.} Unfold \leanname{online_pac_gold_separation}. Show this is the frontier that names the figure's solid content exactly: it is the conjunction $(\exists\,\mathrm{PAC}\wedge\neg\,\text{Online})\wedge(\exists\,\Ex\wedge\neg\,\mathrm{PAC})$, that is \leanname{pac_not_implies_online} and \leanname{ex_not_implies_pac} and \emph{nothing else}, so any sixth ``mutual'' edge the eye expects from a triangle is not in the theorem and must be drawn dashed or not at all.

\item \textbf{Gold's incomparability is a different axis.} Ask why Gold escapes PAC ``upward'' (unbounded VC) yet is escaped by PAC ``downward'' (thresholds). Show this is the boundary that separates two axes: $\Ex\not\Rightarrow\mathrm{PAC}$ is a statistical-capacity fact certified in the kernel (\leanname{ex_not_implies_pac}), while $\mathrm{PAC}\not\Rightarrow\Ex$ is a convergence-signal fact that is literature only (\cite{Gold1967}), so the two legs at the Gold node measure capacity and convergence respectively, not one symmetric ``incomparability''.

\item \textbf{Why three of the legs are dashed.} Ask the kernel for $\mathrm{PAC}\not\Rightarrow\Ex$, $\text{Online}\not\Rightarrow\Ex$, and $\Ex\not\Rightarrow\text{Online}$. Show this is the figure's principal frontier: the paradigms are typed (\leanname{PACLearnable}, \leanname{OnlineLearnable}, \leanname{EXLearnable}) but no theorem in the kernel separates them in these directions, so all three edges are drawn dashed and cited (\cite{Gold1967} for the first two, \cite{Angluin1980} for the third); contrast the solid legs of Exercises~3 and~4 to read line style as evidence grade.

\item \textbf{Pattern languages are the reverse witness.} Toggle the Gold-versus-Online comparison to the other direction. Show this is the boundary that fills the $\Ex\not\Rightarrow\text{Online}$ leg: the class of pattern languages is $\Ex$-identifiable from positive data yet has infinite Littlestone dimension (\cite{Angluin1980}), so it is a class Gold learns that no online learner can, the mirror of the singletons-with-empty-set class that online learns and Gold cannot.

\item \textbf{Singletons sit inside PAC but outside Gold.} Locate the singletons-with-empty-set class on the corrected picture. Show this is the instance that tests the geometry: the class has $\Ldim=1$ so it is online-learnable, hence (by \leanname{online_imp_pac}) it lies in both the Online and PAC regions, yet it is not $\Ex$-identifiable from text (\cite{Gold1967}), so it is a single class witnessing membership in the nested $\{\text{Online}\subset\mathrm{PAC}\}$ region together with exclusion from the Gold region.

\item \textbf{Is the order one graded measure?} Ask whether the VC dimension, the Littlestone dimension, and the mind-change ordinal are three readings of a single parameterized complexity measure under one data-model parameter. Show this is the figure's most general open frontier: the kernel proves the comparabilities it can ($\VC\le\Ldim$ via \leanname{BranchWiseLittlestoneDim_ge_VCDim}, $\text{Online}\Rightarrow\mathrm{PAC}$ via \leanname{online_imp_pac}) but carries the three invariants as separate definitions with no graded predicate unifying them (\cref{op:meta-measure}, \cref{op:effective-three-paradigm}), so the meta-edge collapsing the order into one measure is dashed over an otherwise typed lattice.
\end{enumerate}
